\section{Sphericalization, braided Morita obstructions, and module restrictions}
\label{sec:m3-research-20261022}

This section removes the pivotality hypothesis from four previously unresolved
rank-eight reductions and develops restrictions on possible Morita
constructions.  A spherical structure on the original category is
distinguished throughout from a pseudo-unitary structure obtained after
Galois conjugation.  None of the necessary module actions below supplies
module associators or a new categorification.

\subsection{A signed-quotient criterion for spherical structures}
\label{subsec:research-signed-section}

\begin{lemma}[Spherical section]
\label{lem:research-spherical-section}
Let $\mathcal C$ be a complex fusion category with fusion tensor $N$.
Choose dual-compatible lifts $X_i^+$ of its simple objects in the canonical
sphericalization $\widetilde{\mathcal C}$, and let $T$ be the signed
quotient by the central sign object $\varepsilon=-1$.  If $T$ is
dual-compatible sign-gauge equivalent to $N$, then $\mathcal C$ admits a
spherical structure.
\end{lemma}
\begin{proof}
The sphericalization has two lifts $X_i^\pm$, with
$X_i^-=\varepsilon X_i^+$ and $\dim(\varepsilon)=-1$; see
\cite[Section~5.1]{OstrikPiv3}.  Dual-compatible lifts always exist.
For a nonselfdual pair choose one lift and its dual.  For a selfdual
simple, the dual of a lift lies over the same source object and cannot
be the opposite lift: sphericality gives $\dim(Y^*)=\dim(Y)$, whereas
$\dim(\varepsilon Y)=-\dim(Y)$, and a simple spherical dimension is
nonzero.  This also proves the assertion in
\cite[Section~5.1(a)]{OstrikPiv3} that both lifts are selfdual.

Write $a_{ij}^k$ and $b_{ij}^k$ for the multiplicities of $X_k^+$ and
$X_k^-$ in $X_i^+X_j^+$.  Then
\[
 a_{ij}^k+b_{ij}^k=N_{ij}^k,\qquad
 T_{ij}^k=a_{ij}^k-b_{ij}^k.
\]
Changing lifts by signs $g_i$, with $g_0=1$ and $g_i=g_{i^*}$,
changes $T_{ij}^k$ to $g_i g_j g_kT_{ij}^k$.
After the stated gauge change, $T=N$ and every $b_{ij}^k$ vanishes.
The chosen lifts therefore form a spherical fusion subcategory.  The
forgetful functor restricts to a tensor equivalence from this subcategory
to $\mathcal C$, proving the claim; compare
\cite[Section~5.1(c)]{OstrikPiv3}.
\end{proof}

For use in a complete finite search, the signed tensor necessarily
satisfies
\begin{equation}\label{eq:research-full-parity}
 |T_{ij}^k|\leq N_{ij}^k,\qquad
 T_{ij}^k\equiv N_{ij}^k\pmod2,
\end{equation}
together with unit, duality, cyclic trace symmetry, dual reversal and
all ordered associativity identities.  Commutativity of $T$ must not be
assumed from commutativity of $N$.

\begin{lemma}[Two valid normalizations]
\label{lem:research-two-normalizations}
The following reductions preserve every categorically possible signed
quotient.
\begin{enumerate}
\item In the full domain~\eqref{eq:research-full-parity}, select signed
variables with odd original coefficient whose sign characters are
linearly independent on the group of dual-compatible basis signs.
One may require all selected variables to be positive, retaining at
least one representative of every sign-gauge class.
\item For a selfdual source simple one may choose its lift to have
second Frobenius--Schur indicator $+1$.  With this choice for every
selfdual simple, one has
\begin{equation}\label{eq:research-fs-mod-four}
 T_{j,j^*}^{i}\equiv N_{j,j^*}^{i}\pmod4
 \quad\text{whenever }i=i^*.
\end{equation}
\end{enumerate}
The first normalization is not asserted to preserve the second one.
\end{lemma}
\begin{proof}
For the first assertion, an odd signed variable cannot vanish.
Its sign transforms by a character of the elementary abelian group
of dual-compatible signs.  Independence of the selected characters
makes their joint evaluation map surjective onto all sign patterns.
An appropriate basis-sign change thus makes the selected variables
positive simultaneously.

For the second assertion, both lifts of a selfdual object are selfdual
by Lemma~\ref{lem:research-spherical-section}.  Their pivotal
isomorphisms differ by a sign.  The defining second-indicator rotation
uses the inverse pivotal isomorphism once, so changing this sign
changes the indicator's sign; see \cite[Definition~2.14 and
Lemma~2.16]{LPRV}.  A selfdual simple has indicator $\pm1$, so precisely
one lift has indicator $+1$.  Its opposite has indicator $-1$.
The Fuchs--Runkel--Schweigert parity theorem,
\cite[Theorem~3.9]{LPRV}, says that a selfdual simple of indicator $-1$
occurs with even multiplicity in $Y^*Y$ in a pivotal fusion category.
Applying this inside $\widetilde{\mathcal C}$ gives
\[
 b_{j,j^*}^{i}
 =\frac{N_{j,j^*}^{i}-T_{j,j^*}^{i}}2\in2\Z.
\]
Replacing $j$ by $j^*$ reconciles the two ordered products without a
commutativity assumption.  This proves~\eqref{eq:research-fs-mod-four}.
An arbitrary subsequent change of lift can reverse the indicator,
which explains the final qualification.
\end{proof}

\begin{theorem}[Four general sphericalization reductions]
\label{thm:research-four-general-spherical}
Every complex fusion categorification of any of
\[
 \sha{R8_3_982},\quad \sha{R8_3_1768},\quad
 \sha{R8_3_1795},\quad \sha{R8_3_3472}
\]
admits a spherical structure and has a pseudo-unitary Galois conjugate.
\end{theorem}
\begin{proof}
We enumerate all signed variables modulo cyclic trace symmetry and
dual reversal, imposing every ordered associativity equation.
For the first three rings we use the full parity domain and only
part~(1) of Lemma~\ref{lem:research-two-normalizations}.  The resulting
complete counts are
\[
\begin{array}{c|r|r|r|r|r}
\text{ring}&\text{variables}&\text{sign fixes}&\text{normalized}
 &\text{unrestricted}&\text{gauge classes}\\\hline
R8\_3\_982&42&5&3&96&3\\
R8\_3\_1768&59&6&49&3136&28\\
R8\_3\_1795&59&6&73&4672&42
\end{array}
\]
Surjectivity of the sign normalization is checked against the full
dual-compatible sign group.  Applying that entire group to the
normalized solutions constructs the unrestricted sets explicitly;
the displayed counts do not assume a free group action.
Each row has exactly one gauge class of the original tensor.

To exclude every other class, consider its complete dimension
characters.  A possible sphericalization character $d$ must satisfy
$d_i\ne0$ and $d_i=d_{i^*}$.  Set
$\delta=\sum_i d_i d_{i^*}$.  The two central sign sectors of
$K_0(\widetilde{\mathcal C})\otimes\C$ are the original algebra and
the signed algebra.  A formal codegree $f$ from either sector becomes
$2f$ in the sphericalization, whose dimension is $2\delta$.
Thus central induction requires $\delta/f$ to be an algebraic integer
for every formal codegree of both sectors, by
\cite[Theorem~2.13 and Corollary~2.14]{OstrikPiv3}.

For completeness, the exact norm test used here is the following.
If $p,q$ are the monic minimal polynomials of $\delta,f$, of degrees
$a,b$, integrality requires
\begin{equation}\label{eq:research-codegree-norm}
 b\,v_\ell(|p(0)|)\geq a\,v_\ell(|q(0)|)
 \quad\text{for every rational prime }\ell.
\end{equation}
Indeed, take norms in a common number field and cancel its degree.
Every nongauge signed tensor in these three searches has a commutative
semisimple algebra, verified after enumeration.  A separating regular
matrix with squarefree characteristic polynomial produces every
character orbit, and all character multiplication identities are
checked exactly.  The regular Casimir factors provide all formal
codegrees.  Every character is excluded either by a zero simple
dimension, by unequal dimensions on a dual pair, or by a prime violating
\eqref{eq:research-codegree-norm}.  The receipt records all character
coordinates, global-dimension polynomials and norm witnesses.

For $\sha{R8_3_3472}$ all simples are selfdual.  Here we use precisely
the indicator normalization in part~(2) of
Lemma~\ref{lem:research-two-normalizations}, with no further
positive-variable sign gauge.  The complete search has $84$ signed
variables, visits $298\,865\,973$ native nodes, and returns exactly
one tensor, namely $N$.  In particular no character exclusion is
needed for this signed search.  The independent check reconstructs
the modulo-four domains and verifies the full ordered tensor identities.
An earlier unrestricted search timed out; that search supplies no
conclusion and is not used here.

In all four cases Lemma~\ref{lem:research-spherical-section} gives a
spherical structure on the original category.  The complete ordinary
character test leaves exactly the FP Galois orbit.  More explicitly,
for $\sha{R8_3_982}$ the four characters agreeing on duals form that
orbit; for $\sha{R8_3_1768}$ and $\sha{R8_3_3472}$ every other
character has a zero simple dimension.  For $\sha{R8_3_1795}$ the
other possible dimensions are $20,10,4$: the dimension-$4$ characters
have zero entries, and for the others the quotient by the FP codegree
has norm respectively $1/9$ or $1/36$.  These statements are checked
independently in the full exact character calculation.
The spherical dimension character is therefore a Galois conjugate
of the FP character.  A number-field model and Galois conjugation
\cite[Theorem~2.30 and the following remark]{ENO} give a spherical
category with positive FP dimensions, hence a pseudo-unitary category.
\end{proof}

The source-bound complete replay is
\path{verify_general_spherical_gates.py} in
\path{evidence/rank8_research_20261022/pivotality/};
its receipt is \path{general_spherical_gates.json} and its detailed
proof is \path{PROOF.md} in the same directory.
The native generator substitutes no singleton-domain variables and
uses no domain-dependent sign simplifications.  Domains affect only
the variable order, the enumeration loops, and an absolute-value
overflow bound; the verifier replaces every loop and checks that
the bound does not increase.  Independent tensor reconstruction,
sign-gauge expansion and an optimized full replay are recorded in
\path{evidence/rank8_research_20261022/r928_r941/PIVOTALITY_REVIEW.md}.

\subsection{General braided exclusions and braided Morita representatives}
\label{subsec:research-braided-morita}

\begin{proposition}[Minimal tensor image]
\label{prop:research-minimal-tensor-image}
Let $\mathcal C$ be a fusion category whose global FP dimension is
nonintegral, and suppose every proper fusion subcategory of $\mathcal C$
is integral.  If $\mathcal C$ is categorically Morita equivalent to a
braided fusion category, then $\mathcal C$ itself admits a braiding.
\end{proposition}
\begin{proof}
Let $\mathcal B$ be a braided Morita representative.  Its braiding
defines a fully faithful tensor functor $\mathcal B\to Z(\mathcal B)$.
Compose with $Z(\mathcal B)\simeq Z(\mathcal C)$ and the forgetful
functor to obtain a tensor functor $F:\mathcal B\to\mathcal C$.
Let $\mathcal E$ be its tensor image, closed under direct summands.
If $\mathcal E$ were proper, its integrality and preservation of FP
dimension by tensor functors would give integer FP dimension for
every simple of $\mathcal B$.  Thus $\FPdim(\mathcal B)$ would be
an integer, contradicting Morita invariance
$\FPdim(\mathcal B)=\FPdim(\mathcal C)$.
Hence $F$ is dominant.  A dominant tensor functor between fusion
categories of equal FP dimension is an equivalence.  Transport its
braiding to $\mathcal C$.  The cited standard facts are
\cite[Proposition~6.3.4, Corollary~7.16.7 and
Theorem~8.12.3]{EGNO}.
\end{proof}

\begin{theorem}[No braided Morita representatives]
\label{thm:research-no-braided-morita}
No complex categorification of any of
\[
\begin{gathered}
\sha{R8_3_65},\ \sha{R8_3_928},\ \sha{R8_3_941},\ \sha{R8_3_982},\\
\sha{R8_3_1154},\ \sha{R8_3_1768},\ \sha{R8_3_1795},\ \sha{R8_3_3472}
\end{gathered}
\]
is Morita equivalent to a braided fusion category.
\end{theorem}
\begin{proof}
First exclude a braiding on the source itself.
The rings $\sha{R8_3_928}$, $\sha{R8_3_941}$ and
$\sha{R8_3_1154}$ are noncommutative, immediately excluding braidings.
The ring $\sha{R8_3_65}$ is an irrational generalized near-group
ring: its invertibles form $C_4$ and act transitively on four
noninvertibles of dimension $(11+5\sqrt5)/2$.  Neither possibility
in \cite[Theorem~1]{SchopierayBraidedGNG}, whose noninvertible
dimensions are $(1+\sqrt5)/2$ and $1+\sqrt2$, has this dimension.
That theorem does not assume pivotality.

For the other four rings, Theorem~\ref{thm:research-four-general-spherical}
allows reduction to a pseudo-unitary Galois conjugate.  A braiding
transports under the same field automorphism because its naturality
and hexagon identities are algebraic.  It suffices to exclude
pseudo-unitary braided categorifications.

For $\sha{R8_3_1768}$ the simple dimension $5+\sqrt{37}$ has
minimal polynomial $x^2-10x-12$, and for $\sha{R8_3_3472}$ the
simple dimension $6+5\sqrt2$ has polynomial $x^2-12x-14$.
Squared norms of simples in a braided fusion category are d-numbers
\cite[Theorem~1.8]{OstrikCode}.  In the pseudo-unitary setting they
are squared FP dimensions; an algebraic integer is a d-number if and
only if its square is, by \cite[Corollary~2.9]{OstrikCode}.
For a quadratic polynomial $x^2-tx+n$, the d-number criterion
\cite[Lemma~2.7(v)]{OstrikCode} requires $n\mid t^2$.
Here $12\nmid100$ and $14\nmid144$, respectively.

For $\sha{R8_3_982}$, the complete earlier induction-frame and
indicator analysis forces a unique central invertible $G$ forgetting
to the source invertible $g=X_6$, with $G^2=\one$ and twist $\pm i$.
The completeness is essential: the two excluded frames have exact
finite-domain certificates, while both surviving frames force the
same central lift and the indicator $\nu_2(g)=-1$.
A braiding on the source supplies a central lift of $g$, necessarily
$G$.  Its pointed $C_2$ subcategory is therefore semionic and
nondegenerate.  The nondegenerate-subcategory decomposition
\cite[Theorem~3.13]{DGNO} would give
$\mathcal C\simeq\mathrm{Semion}\boxtimes\mathcal D$.
Its adjoint subcategory is proper, being contained in
$\Vec\boxtimes\mathcal D$.  In the canonical ring, closure of the
supports of every $XX^*$ instead gives the full ring.  This is a
contradiction.  The complete frame and lift certificates are
\path{evidence/rank8_next_20261021/r982/FRAME01_RESULT.md} and
\path{SEMION_ACTION.md} in that directory.

It remains to treat $\sha{R8_3_1795}$.  Its adjoint subring is the
whole ring, and its only proper nontrivial based subring is
$\{\one,g\}$, where $g=X_1$ has order two.  In a braided category
invertibles centralize the adjoint category, so $g$ is transparent.
The whole category is not symmetric, since its FP dimensions are
nonintegral.  Thus its M\"uger center is precisely $\Rep(C_2)$ or
$\mathrm{sVec}$.
Set $D=540+240\sqrt5$.  Multiplication by $g$ pairs the simples as
$(0,1),(2,7),(3,6),(4,5)$.  The four characters with $\chi(g)=1$
have formal codegrees $D,\overline D,10,10$.  The pair quotient
halves those codegrees.

If the center is Tannakian, de-equivariantization is nondegenerate.
Each pair gives one simple because the free module on $X$ has
endomorphism space
$\operatorname{Hom}(X,(\one\oplus g)X)$ of dimension one.
The modularization has rank four, dimension $D/2$, and simple
dimensions $d_i$ for one representative of each pair.  The positive
spherical structure descends; see \cite[Proposition~4.56(iii)
and Section~6.3]{DGNO}.  The modular codegree identity is therefore
$(D/2)/d_i^2$.
If the center is $\mathrm{sVec}$, the reduced S-matrix on one
representative per pair has square $(D/2)I$ in this positive,
selfdual situation.  Its normalized rows give all quotient
characters, and the same row-norm argument gives the same codegrees;
see \cite[Lemma~3.1 and Proposition~3.3]{LacabanneSlightly}.
Thus in either case a full-ring codegree $10$ forces
$d_i^2=D/10=54+24\sqrt5$.  The actual squared dimensions are
$1,81+36\sqrt5,94+42\sqrt5$, none of which has that value.
This excludes the final braided case.

For all eight rings, complete enumeration of unit-containing
dual-closed subsets shows that every proper based subring is pointed.
The global FP dimension is irrational, as checked by isolating the
largest regular Casimir eigenvalue in an irreducible factor of degree
greater than one.  Proposition~\ref{prop:research-minimal-tensor-image}
now excludes every braided Morita representative.
\end{proof}

The braided numerical and character calculations are replayed by
\path{site/verify_braided_scope.py}.  The last subring and dimension
checks have the independent source-bound replay
\path{evidence/rank8_research_20261022/braided_morita/verify.py}.
The theorem concerns braided Morita representatives; it does not
exclude a Morita representative which is itself unbraided.

\subsection{The two equal-dimension rings have different centers}
\label{subsec:research-1154-1795-morita}

\begin{proposition}
\label{prop:research-1154-1795-no-morita}
No complex categorifications of $\sha{R8_3_1154}$ and
$\sha{R8_3_1795}$ are categorically Morita equivalent.
\end{proposition}
\begin{proof}
Both have FP dimension $D=540+240\sqrt5$.  In the pseudo-unitary
case, the eight scalar characters of $\sha{R8_3_1795}$ have formal
codegrees
\[
 4,\ 20,\ 10,10,10,10,\ D,\overline D.
\]
The injective assignment in the induced-unit theorem
\cite[Theorem~2.13]{OstrikPiv3} gives four \emph{distinct} central
simples of dimension $D/10=54+24\sqrt5$.
For $\sha{R8_3_1154}$, the exhaustive restrictions of
Section~\ref{sec:r1154-three-center-patterns} leave the following
complete patterns:
\[
\begin{array}{c|r|r}
\text{pattern}&\text{center rank}&
 \#\{Z:\FPdim(Z)=54+24\sqrt5\}\\\hline
6971&23&0\\8766&21&0\\8772&29&0.
\end{array}
\]
The unique allowed dimension type of this size is high-column
type $2$, whose multiplicity vanishes in all three patterns.
Equivalent centers preserve FP dimensions, so these inventories
cannot be equivalent.

For general complex categorifications, apply the established
Galois/pseudo-unitary reduction for $\sha{R8_3_1154}$ to the entire
putative Morita context.  It preserves both fusion rings and Morita
equivalence.  Global dimension and FP dimension are Morita invariants
\cite[Theorem~2.15 and Corollary~8.14]{ENO}; the transformed partner
is consequently pseudo-unitary as well.  The preceding contradiction
applies.  This proof does not require the new separate sphericalization
result for $\sha{R8_3_1795}$.
\end{proof}

The direct exact comparison is
\path{evidence/rank8_research_20261022/r1154_r1795/morita_obstruction.py}.
It verifies every character identity against the source tensor and
every displayed dimension multiplicity, while retaining the earlier
complete center enumerations as explicit dependencies.  This
proposition rules out a shared Morita construction for the pair,
without excluding either ring individually.

\subsection{All module ranks and based module types for
\texorpdfstring{$\sha{R8_3_982}$}{R8-3-982}}
\label{subsec:research-r982-module-types}

\begin{theorem}
\label{thm:research-r982-modules}
For any complex categorification $\mathcal C$ of $\sha{R8_3_982}$,
every indecomposable semisimple module category has rank eight.
Its based module is the regular one or one of two explicitly recorded,
mutually dual, nonregular types.  The complete necessary list of
underlying connected separable algebra objects contains eight
objects.  Every nonregular module type has a simple object with
internal end $\one\oplus X_5$, of FP dimension
$3+\sqrt2+\sqrt3+\sqrt6$.
\end{theorem}
\begin{proof}
The preceding central semion analysis gives a nontrivial pointed
associator on $\langle g\rangle$, where $g=X_6$.  This applies first
in the pseudo-unitary Galois conjugate supplied by
Theorem~\ref{thm:research-four-general-spherical}.  The two classes
in $H^3(C_2,\C^\times)$ have representatives valued in $\{\pm1\}$,
so nontriviality is Galois invariant and applies to $\mathcal C$.
The element $g$ acts freely on the simple objects of every module
category.  Indeed a fixed module simple would be a rank-one module
for the pointed $C_2$ category, and its scalar module pentagon would
trivialize that associator.  Equivalently, use the pointed-module
classification in \cite[Example~2.1]{OstrikDouble}.

For a simple module object $m$, let $A=\underline{\operatorname{End}}(m)$.
Its coefficient at $g$ is $\dim\operatorname{Hom}(gm,m)=0$.
The earlier complete $31$-object algebra filter therefore leaves
exactly eleven objects.  The source permutation of $g$ is
$(0\ 6)(1\ 2)(3\ 4)(5\ 7)$.
Fix one of these candidates $A=\sum_i a_iX_i$, and define the
induction matrix
\[
 B_{i,n}=\dim\operatorname{Hom}(X_i m,n).
\]
Adjunction gives
\begin{equation}\label{eq:research-r982-gram}
 (BB^{\mathsf T})_{ij}=\sum_k a_kN_{j,k}^{i}.
\end{equation}
Every module simple occurs in some $X_i m$, so every column is
nonzero.  The columns at $m,gm$ are $a$ and its $g$-translate.
All others have zero coordinates at rows $0,6$ and occur in pairs
$(v,gv)$.

Remove the distinguished pair and order the six remaining rows as
$(1,2;3,4;5,7)$.  For one representative of each residual column
pair, let $u$ and $w$ be its three within-pair sums and differences.
If $H$ is the residual matrix in~\eqref{eq:research-r982-gram}, set
\[
 P_{ab}=H_{2a,2b}+H_{2a,2b+1}=\sum u_a u_b,
 \qquad
 M_{ab}=H_{2a,2b}-H_{2a,2b+1}=\sum w_a w_b.
\]
The column conditions are exactly
$u\geq0$, $|w_a|\leq u_a$, and $u_a\equiv w_a\pmod2$;
they reconstruct the nonnegative columns $(u+w)/2,(u-w)/2$.
Each $u$ is nonzero.

Seven candidates have $\det P=1$.  By Cauchy--Binet the sum of
squares of all $3\times3$ minors of an integral factor $U$ is one.
Precisely one minor is nonzero.  Its three columns form an integral
basis.  Replacing any of them by another column has determinant zero,
so that additional column must be zero.  Hence there are precisely
three residual column pairs, and the module rank is eight.

The other four candidates have $\det P=4$ and checked unimodular
isometries to $\operatorname{diag}(1,2,2)$.  Any integral embedding
of this lattice, up to signed coordinate permutation, has a norm-one
row on one coordinate and two orthogonal norm-two rows outside it.
The latter each have two entries $\pm1$; their supports either
coincide or are disjoint.  These are the only possibilities, giving
three or five nonzero columns.  Transforming back gives complete
factors $U$.  In every five-column case its second row has five odd
entries, whereas $M_{11}=1$ in zero-based indexing.  The parity
condition would force five nonzero entries in the corresponding row
of $w$, contradicting its squared norm one.  Only three-column
factors remain, again giving module rank eight.

In each case $\det M=\det P\in\{1,4\}$, so the residual Gram
matrix has rank six.  The distinguished columns have independent
coordinates at rows $0,6$, giving full Gram rank eight.
Thus $B$ is square and invertible.  Each source action is forced by
$B^{-1}N_iB$.  Exhaust all integral three-column factors of $P$ and
$M$ within their exact squared-norm bounds, and require these actions
to be nonnegative integral.  The complete results, indexed by the
prior algebra list, are
\[
\begin{array}{c|l|l}
\text{index}&(a_0,\ldots,a_7)&\text{outcome}\\\hline
0 &(1,0,0,0,0,0,0,0)&\text{regular}\\
2 &(1,0,0,0,0,1,0,0)&\text{two nonregular types}\\
5 &(1,0,0,1,1,1,0,0)&\text{two nonregular types}\\
9 &(1,1,1,1,1,1,0,0)&\text{two nonregular types}\\
10&(1,1,1,1,1,1,0,1)&\text{forced action entry }-1\\
12&(1,1,1,1,1,2,0,0)&\text{regular}\\
14&(1,1,1,2,2,1,0,1)&\text{forced action entry }-1\\
16&(1,1,1,2,2,2,0,0)&\text{regular}\\
19&(1,2,2,3,3,2,0,2)&\text{regular}\\
21&(1,2,2,3,3,3,0,1)&\text{no compatible factor}\\
23&(1,3,3,4,4,3,0,2)&\text{two nonregular types}.
\end{array}
\]
All surviving actions satisfy every source multiplication, dual
transpose, connectivity and free-$g$ identity.  Exact simultaneous
permutation checks give exactly three types, the two nonregular
ones interchanged by source duality.  Their internal-end vectors
are exactly the four at indices $2,5,9,23$, each twice.  Thus both
contain $\one\oplus X_5$.  A module category of regular type has
a simple with internal end $\one$, and is equivalent to the regular
module category.  This proves all assertions.
\end{proof}

The complete small-lattice factors, forced negative entries, full
actions and intertwining permutations are recorded and replayed by
\path{evidence/rank8_research_20261022/r982/module_types.py}.
The proof depends on the previous exhaustive algebra filter through
its source-bound receipt, and on no chosen bound for module rank.
Whether $\one\oplus X_5$ admits the required multiplication and
associativity remains open; the theorem reduces every nonregular
Morita route to this concrete algebra problem.

\subsection{Rank-two Morita-dual spectra for
\texorpdfstring{$\sha{R8_3_3472}$}{R8-3-3472}}
\label{subsec:research-r3472-dual-spectra}

\begin{proposition}
\label{prop:research-r3472-ten-spectra}
If the recorded rank-two based module of $\sha{R8_3_3472}$ whose
$X_1$ action is $\left(\begin{smallmatrix}5&7\\7&7\end{smallmatrix}\right)$
is realized by a module category, its Morita dual
has one of the ten FP dimension spectra in the table below.
This is a necessary condition, without an existence assertion for
either module category or any spectrum.
\end{proposition}
\begin{proof}
Put $r=\sqrt2$, $d=7+5r$ and $D=700+490r$.
After reversing the two module simples, the source action of $X_1$
is
\[
 \begin{pmatrix}7&7\\7&5\end{pmatrix}.
\]
An action matrix of a simple $Y$ in the dual category commutes with
this matrix.  Solving the commutator equations over $\Z$ gives
\[
 \begin{pmatrix}s+t&7t\\7t&s-t\end{pmatrix},
\]
where $t\geq0$ and $s\geq t$.
The common positive eigenline is the module FP-dimension line,
and hence $\FPdim(Y)=s+5tr$.

Morita invariance gives $\FPdim(\mathcal B)=D$ for the dual.
The regular FP-dimension virtual object acts with rank one on an
indecomposable module; see \cite[Exercise~7.16.8]{EGNO}.
Thus its eigenvalue on the other line vanishes, giving
\[
 \sum_Y(s_Y+5t_Yr)^2=700+490r,
 \qquad \sum_Y(s_Y+5t_Yr)(s_Y-5t_Yr)=0.
\]
Equivalently,
\[
 \sum s_Y^2=350,\qquad\sum t_Y^2=7,
 \qquad\sum s_Yt_Y=49.
\]
Writing $k_Y=s_Y-7t_Y$ reduces these to
$\sum k_Y^2=7$ and $\sum t_Yk_Y=0$.

The integral-dimensional simples, exactly those with $t=0$, form
a fusion subcategory because all radical coefficients are
nonnegative.  If its dimension is $h$, the last square sum gives
$1\leq h\leq7$.  Algebraic-integrality of $D/h$
\cite[Theorem~7.17.6]{EGNO} requires integer trace $1400/h$
and norm $9800/h^2$, leaving $h\in\{1,2,5,7\}$.
This subcategory is pointed.  Only the cases $h=5,7$ need comment:
a nonpointed dimension-five category would have dimensions $1,2$,
contradicting $4=1+2n$ for the square of its noninvertible.
For dimension seven the only nonpointed possibility is $1,1,1,2$.
Its three invertibles fix the dimension-two simple, whose product
with its dual would require $4=3+2n$, again impossible.

Let $g=h$ be the number of invertibles.  The nonintegral dimensions
are $td+k$, and their multiplicities $n_{t,k}$ satisfy the complete
finite conditions
\[
 t\in\{1,2\},\quad k\in\{-2,-1,0,1,2\},\quad
 \sum n_{t,k}t^2=7,\quad
 \sum n_{t,k}k^2=7-g,\quad
 \sum n_{t,k}tk=0.
\]
Exact enumeration gives eleven spectra, all with nonnegative
matrices in this commutant cone.  For $g=7$ the two are
$1^7,d^7$ and $1^7,d^3,(2d)^1$.  Since $d^2=1+14d$, the product
$XX^*$ for a $d$-dimensional simple contains exactly one invertible:
comparison of the constant terms in $\Q(d)$ proves this assertion.
Its stabilizer in the order-seven invertible group is therefore
trivial.  The number of $d$-dimensional simples is divisible by
seven, excluding the second spectrum.
\end{proof}

In the following table exponents denote multiplicities, and the
nonintegral part is written using $d=7+5\sqrt2$.
\[
\begin{array}{c|r|l}
\text{invertibles}&\text{rank}&\text{nonintegral dimensions}\\\hline
1&5&(d-1)^2,\ d+2,\ 2d\\
1&8&(d-1)^2,\ d^4,\ d+2\\
1&8&(d-1)^3,\ d,\ (d+1)^3\\
1&5&d-2,\ (d+1)^2,\ 2d\\
1&8&d-2,\ d^4,\ (d+1)^2\\
2&6&d^2,\ d+2,\ 2d-1\\
2&6&d-2,\ d^2,\ 2d+1\\
5&9&d-1,\ d,\ d+1,\ 2d\\
5&12&d-1,\ d^5,\ d+1\\
7&14&d^7
\end{array}
\]
The notation in this table refers to lists, not powers of algebraic
numbers.  The replay
\path{evidence/rank8_research_20261022/r1768_r3472/r3472_rank_two_dual_spectra.py}
checks the stated source action matrix, the complete square-sum
enumeration, matrix nonnegativity and the orbit exclusion.
Repeated simple objects with identical action matrices are allowed;
no commutativity or faithfulness assumption on the dual fusion ring
is used.  The other recorded rank-two action has a larger integral
commutant and is outside the scope of this proposition.

\subsection{Pointed induction and an exclusion of
\texorpdfstring{$\sha{R8_3_65}$}{R8-3-65}}
\label{subsec:research-r65-exclusion}

The next argument covers every pointed associator class directly.
In particular it does not use the earlier class-zero subgroup-dual
argument or its lower-rank and $\sha{R8_3_1453}$ dependencies.

\begin{lemma}[Pointed induction and integral Fourier fibers]
\label{lem:research-pointed-fourier}
Let $\mathcal C$ be pseudo-unitary with its positive spherical
structure, let $D=\FPdim(\mathcal C)$, and let $g$ be an invertible
simple of order $n$.  Write
\[
 I(g)=\bigoplus_Z m_Z Z,\qquad z_Z=\FPdim(Z)>0,
\]
where $I$ is adjoint to the central forgetful functor.  Then every
constituent satisfies $\theta_Z^n=\nu_n(g)$.  For any choice
$\alpha^n=\nu_n(g)$, each twist fiber has weight
\begin{equation}\label{eq:research-fourier-weight}
 \sum_{\theta_Z=\alpha e^{2\pi ia/n}}m_Zz_Z=D/n
 \qquad(a\in\Z/n).
\end{equation}
Suppose further that all noninvertible source simples have the same
irrational FP dimension $d$, and inductions of distinct invertibles
have disjoint simple supports.  Write $D/n=A+Bd$ with $A,B\in\Q$.
If $F(Z)=m_Zg+\sum_X b_{Z,X}X$ and $b_Z=\sum_Xb_{Z,X}$,
then in each fiber
\begin{equation}\label{eq:research-integral-fibers}
 \sum m_Z^2=A,\qquad \sum m_Zb_Z=B.
\end{equation}
\end{lemma}
\begin{proof}
The pointed restriction gives $|\nu_n(g)|=1$ and
$\nu_j(g)=0$ for $1\leq j<n$.  Central induction gives
\[
 \sum_Zm_Zz_Z\theta_Z^j=D\nu_j(g),\qquad
 \sum_Zm_Zz_Z=D;
\]
see \cite[Theorem~4.1 and the calculation preceding Lemma~9.3]{NgIndicators}.
At $j=n$, equality in the triangle inequality, with every nonzero
weight positive, forces $\theta_Z^n=\nu_n(g)$ for each constituent.
The $n$ possible twists are therefore $\alpha e^{2\pi ia/n}$.
The first $n-1$ indicator identities and the dimension identity form
the discrete Fourier transform of their fiber weights.  Its inverse
gives~\eqref{eq:research-fourier-weight}.

By adjunction, a constituent of $I(g)$ contains $g$ in its forgotten
image with multiplicity $m_Z$.  Disjointness of the invertible
inductions excludes every other invertible from that image.  Thus
$z_Z=m_Z+b_Zd$ with $b_Z\in\NN$.
Comparing coefficients of $1,d$ in
\eqref{eq:research-fourier-weight} proves
\eqref{eq:research-integral-fibers}.  No bound on twist orders or
unitarity assumption is used.
\end{proof}

\begin{theorem}
\label{thm:research-r65-excluded}
The fusion ring $\sha{R8_3_65}$ has no categorification by a complex
fusion category.
\end{theorem}
\begin{proof}
Its source invertibles form $C_4$, generated by $g=X_1$ with
$g^2=X_7$ and $g^3=X_2$.  The other four simples, $X_3,X_4,X_5,X_6$,
have dimension
\[
 d=(11+5\sqrt5)/2,\qquad d^2=1+11d,\qquad D=8+44d.
\]
The published two-dimension reduction
\cite[Lemma~3.1]{SchNear} implies that every hypothetical complex
categorification has a pseudo-unitary Galois conjugate.  We work in
that conjugate.

The exact induction Gram matrix has diagonal $8$ at the four
pointed positions and zero between distinct pointed positions.
Thus their inductions have disjoint supports.  For either generator
$g$ or $g^3$, Lemma~\ref{lem:research-pointed-fourier} gives, in
each of its four twist fibers,
\[
 \sum m_Z^2=2,\qquad \sum m_Zb_Z=11.
\]
There are exactly two constituents per fiber, each with $m_Z=1$,
and each satisfies $b_Z\leq11$.  In particular $I(g)$ and $I(g^3)$
are multiplicity free with eight constituents each.  The certificate
below uses only the weaker necessary column restrictions $m_Z=1$
and $b_Z\leq11$.

Let $k\in\Z/4$ be the pointed associator class, with cyclic
representative
\[
 \omega_k(a,b,c)=\exp\left(\frac{2\pi i k a}{4}
                 \left\lfloor\frac{b+c}{4}\right\rfloor\right).
\]
The pointed indicators are
$\nu_4(g)=\nu_4(g^3)=i^k$ and $\nu_2(g^2)=(-1)^k$.
The phase assertion of Lemma~\ref{lem:research-pointed-fourier}
therefore imposes the following exact normalized relative traces
of $(\theta,\theta^2,\theta^3,\theta^4)$ over $K=\Q(\sqrt5)$:
\[
\begin{array}{c|l|l}
k&\text{constituent over }g\text{ or }g^3
 &\text{constituent over }g^2\\\hline
0&(1,1,1,1),\ (-1,1,-1,1),\ (0,-1,0,1)
 &(1,1,1,1),\ (-1,1,-1,1)\\
1,3&(0,0,0,0)&(0,-1,0,1)\\
2&(0,0,0,-1)&(1,1,1,1),\ (-1,1,-1,1).
\end{array}
\]
For example, odd $k$ forces primitive order-$16$ generator twists,
and $k=2$ forces primitive order-$8$ generator twists.
These trace evaluations use the linear disjointness of
$\Q(\sqrt5)$ and $\Q(\zeta_{16})$.  They respect duality, since
central duals have equal twists and $g,g^3$ have equal fourth
indicator.

We impose these necessary restrictions on the established exact
central-induction cone with full-character-field S-symmetry and
the real part of the CM-field STS identities.  It has $283$ rational
equality rows.  Only forbidden pointed columns are removed; all
other center columns and all auxiliary rotation columns remain.
For $k=0,1,3$ we also use the proved Fourier restrictions
$m_Z=1,b_Z\leq11$.  The complete second-indicator possibilities,
in source order $1,\ldots,7$, are
$(0,0,0,0,1,e,f)$ with $e,f\in\{\pm1\}$, and the pointed class
fixes $f=(-1)^k$.  Thus the following six branches are exhaustive:
\[
\begin{array}{c|r|r|l}
k&e&\text{allowed columns}&\text{certificate}\\\hline
0&-1&2413&\texttt{pointed\_phase\_k0\_fourier.json},\ 0\\
0&+1&2413&\texttt{pointed\_phase\_k0\_fourier.json},\ 1\\
1,3&-1&1992&\texttt{pointed\_phase\_k1\_fourier.json},\ 0\\
1,3&+1&1992&\texttt{pointed\_phase\_k1\_fourier.json},\ 1\\
2&-1&2222&\texttt{pointed\_phase\_k2.json},\ 0\\
2&+1&2222&\texttt{pointed\_phase\_k2.json},\ 1.
\end{array}
\]
In every row an exact rational covector $y$ satisfies
$y\cdot c\geq0$ for \emph{every} allowed column $c$, but
$y\cdot b<0$ for that branch's right-hand side.  Hence the
necessary system $\sum_c x_cc=b$, $x_c\geq0$, has no solution.
The certificates check these inequalities both by the factorized
exact oracle and by explicit sparse-column arithmetic.  A numerical
linear program proposes the covectors but contributes no proof
step.  Changed tensors, missing branches or changed domain metadata
fail verification.

We justify the finiteness of the twist-trace domain to distinguish
this argument from a bounded-order search.  The Cauchy theorem for
spherical fusion categories~\cite{BNRWCauchy}, with
$\operatorname{Norm}(D)=2000$, restricts the prime divisors of the
central twist order to $2,5$.  For order $2^a5^b$, the normalized
relative traces at powers $1,2,3,4$ all vanish if $a\geq4$ or
$b\geq2$.  This follows respectively from the remaining primitive
$2$-power trace, or the trace down the $5$-power tower.  The finite
box $a\leq3,b\leq1$ is represented by divisors of $40$.
These traces, together with the zero tuple for all higher supported
orders, give exactly the thirteen trace types used in the cone.
The metadata value $40$ is therefore a compression bound, not a
bound on actual twists.  The order-three rotation eigenvalues on
$\operatorname{Hom}(\one,X^3)$ belong to a different auxiliary
domain; their columns are retained and are not subjected to this
prime-support restriction.

Every $k\in\Z/4$ is excluded.  This contradicts the existence of
the pseudo-unitary Galois conjugate and proves the general complex
exclusion.
\end{proof}

The exact implementation and theoretical dependency record are
\path{evidence/rank8_research_20261022/r65/pointed_phase_cone.py}
and \path{PROOF.md} in the same directory.  The inherited cone is
specified by \path{PROOF.md} and \path{ATOMIC_MODULAR.md} in
\path{evidence/rank8_frontier_20261021/global_induction/},
and by \path{GLOBAL_CONVEX_INDUCTION.md} in
\path{evidence/rank8_m3_20261021/criteria/}.
The earlier feasible witness for the weaker odd-class relaxation
is retained only to show that the Fourier-integrality step is
essential; it is not a categorification or a survivor of the
strengthened conditions.
